% Part: computability
% Chapter: recursive-functions
% Section: non-pr-functions

\documentclass[../../../include/open-logic-section]{subfiles}

\begin{document}

\olfileid[pt]{cmp}{rec}{npr}
\olsection{Funções que não são recursivas primitivas}

As funções recursivas primitivas não esgotam as funções intuitivamente
computáveis. Deve ser intuitivamente claro que podemos fazer uma lista
de todas as funções recursivas primitivas unárias,
$f_0$, $f_1$, $f_2$,~\dots, de modo que possamos calcular efetivamente
o valor de $f_x$ na entrada $y$; em outras palavras, a função $g(x,y)$,
definida por
\[
g(x,y) = f_x(y)
\]
é computável. Mas então também é computável a função
\begin{eqnarray*}
h(x) & = & g(x,x) + 1 \\
& = & f_x(x) +1.
\end{eqnarray*}
Para cada função recursiva primitiva $f_i$, os valores de $h$ e
$f_i$ diferem em $i$. Logo, $h$ é computável, mas não é recursiva
primitiva; e o mesmo se pode dizer de $g$. Esta é uma versão
``efetiva'' do argumento de diagonalização de Cantor.

Podemos dar exemplos mais explícitos de funções computáveis que não
são recursivas primitivas. Por exemplo, seja $g^n(x)$ a notação para
$g(g(\dots g(x)))$, com um total de $n$ ocorrências de $g$; e defina
uma sequência $g_0,g_1,\dots$ de funções por
\begin{eqnarray*}
g_0(x) & = & x+1 \\
g_{n + 1}(x) & = & g_n^x(x)
\end{eqnarray*}
Você pode verificar que cada função $g_n$ é recursiva primitiva.
Cada função sucessiva cresce muito mais depressa do que a anterior;
$g_1(x)$ é igual a $2x$, $g_2(x)$ é igual a $2^x \cdot x$, e
$g_3(x)$ cresce aproximadamente como uma torre exponencial de $x$
algarismos $2$. A função de Ackermann--P\'eter é essencialmente a
função $G(x) = g_x(x)$, e pode-se mostrar que ela cresce mais depressa
do que qualquer função recursiva primitiva.

Voltemos à questão de enumerar as funções recursivas primitivas.
Lembre-se de que atribuímos notações simbólicas a cada função recursiva
primitiva; portanto, basta enumerar as notações. Podemos atribuir um
número natural $\#(F)$ a cada notação $F$, recursivamente, da seguinte
maneira:
\begin{eqnarray*}
\#(0) & = & \langle 0 \rangle \\
\#(S) & = & \langle 1 \rangle \\
\#(\Proj{n}{i}) & = & \langle 2, n, i \rangle \\
\#(\fn{Comp}_{k,l}[H,G_0,\dots,G_{k-1}]) & = & \langle
3,k,l,\#(H),\#(G_0),\dots,\#(G_{k-1}) \rangle \\
\#(\fn{Rec}_l[G,H]) & = & \langle 4, l, \#(G), \#(H) \rangle
\end{eqnarray*}
Aqui usamos o fato de que toda sequência de números pode ser vista como
um número natural, por meio dos códigos da seção anterior. O resultado
é que cada código recebe um número natural. Naturalmente, algumas
sequências (e, portanto, alguns números) não correspondem a notações;
mas podemos tomar $f_i$ como a função recursiva primitiva unária cuja
notação tem código $i$, se $i$ codificar tal notação, e como a função
constante de valor $0$ caso contrário. O resultado final é um modo
explícito de enumerar as funções recursivas primitivas unárias.

(Na verdade, algumas funções, como a função constante zero, aparecerão
mais de uma vez na lista. Isso não é apenas um artefato de nossa
codificação, mas também decorre do fato de que a função constante zero
tem mais de uma notação. Veremos mais adiante que não se pode evitar
computavelmente essas repetições; por exemplo, não existe uma função
computável que decida se uma dada notação representa ou não a função
constante zero.)

Podemos agora tomar a função $g(x,y)$ como sendo dada por $f_x(y)$,
em que $f_x$ se refere à enumeração que acabamos de descrever. Como
sabemos que $g(x,y)$ é computável? Intuitivamente, isso é claro: para
calcular $g(x,y)$, primeiro ``desempacote'' $x$ e verifique se ele é
uma notação para uma função unária. Se for, calcule o valor dessa função
na entrada~$y$.

\begin{tagblock}{TMs}
\begin{digress}
Talvez você já esteja convencido de que (com algum trabalho!) é possível
escrever um programa (digamos, em Java ou C++) que faça isso; e agora
podemos apelar à tese de Church--Turing, segundo a qual tudo o que é
intuitivamente computável pode ser computado por uma máquina de Turing.

Naturalmente, uma maneira mais direta de mostrar que $g(x,y)$ é
computável seria descrever explicitamente uma máquina de Turing que a
calcule. Isso evitaria, em particular, a tese de Church--Turing e os
apelos à intuição. Em breve teremos desenvolvido recursos suficientes
para mostrar que $g(x,y)$ é computável recorrendo a um modelo de
computação que pode ser \emph{simulado} em uma máquina de Turing:
a saber, as funções recursivas.
\end{digress}
\end{tagblock}

\end{document}
